\section*{Introduction}
\phantomsection\addcontentsline{toc}{section}{Introduction}

Diabetic retinopathy is one of the most frequent complications of diabetes and a leading
cause of preventable blindness in adults of working age. Because the disease is silent in
its early stages, it can only be controlled by regular screening of the retina, and the
number of people to screen is now far beyond what human graders can handle. Automated
grading of fundus photographs by deep learning has therefore become a major research
topic, and several systems are already used in practice.

This chapter introduces the clinical problem on which the whole thesis is built. We first
define diabetic retinopathy and describe how it develops. We then present the
international severity scale and the fundus images used to grade it, and the way
screening programmes measure patient safety. We review a decade of deep-learning methods
for automated grading, and we close with the challenges and limits that current systems
still share. These limits define the research problem addressed in the following
chapters.

\section{Diabetic Retinopathy: Definition and Clinical Context}
\label{sec:burden}

\subsection{Anatomy of the Retina}
\label{sec:anatomy}

The retina is the thin layer of neural tissue that lines the back of the eye and converts
light into nerve signals. Seen through the pupil, as in a fundus photograph
(\cref{fig:anatomy}), it shows three main landmarks. The \emph{optic disc} is the point
where the optic nerve leaves the eye and where the central retinal artery and vein enter;
it appears as a bright oval. From it, the retinal vessels branch into superior and
inferior \emph{vascular arcades} that curve around the \emph{macula}, a region about
$5.5$~mm wide responsible for central vision. At the centre of the macula, the
\emph{fovea} contains the highest density of cone photoreceptors and has no inner retinal
vessels, which is why it appears as a darker spot~\cite{abramoff2010retinal}.

The retinal circulation is a terminal network: each capillary bed depends on a single
feeding arteriole, and there are few alternative paths when a vessel closes. This makes
the retina particularly vulnerable to diseases of small vessels such as diabetes, and it
also makes the vascular tree a direct, non-invasive window on the state of the
microcirculation. The shape of this tree -- its branching, its loops and its gaps -- is
therefore clinically meaningful, an observation that motivates the topological features
used later in this thesis.

% Figure 1.1 --- schematic of a colour fundus photograph (left eye view)
\begin{figure}[H]
\centering
\resizebox{0.8\linewidth}{!}{%
\begin{tikzpicture}[font=\footnotesize]
\begin{scope}
  \clip (0,0) circle (3.0);
  \fill[orange!55!red!45!white] (0,0) circle (3.0);
  % vessels: superior and inferior arcades leaving the optic disc
  \foreach \s/\w in {1/1.6pt,-1/1.6pt}{
    \draw[red!55!black, line width=\w] (1.9,0.1*\s) .. controls (1.2,1.3*\s) and (-0.6,1.7*\s) .. (-2.9,1.0*\s);
    \draw[red!45!black, line width=1.1pt] (1.9,0.25*\s) .. controls (1.4,1.7*\s) and (0.2,2.4*\s) .. (-1.5,2.8*\s);
    \draw[red!55!black, line width=0.8pt] (0.3,1.35*\s) .. controls (-0.1,0.9*\s) and (-0.3,0.7*\s) .. (-0.45,0.55*\s);
    \draw[red!55!black, line width=0.8pt] (-1.2,1.55*\s) .. controls (-1.5,1.0*\s) and (-1.9,0.6*\s) .. (-2.4,0.45*\s);
  }
  \draw[red!55!black, line width=1.0pt] (2.0,0) -- (2.9,0.5);
  \draw[red!55!black, line width=1.0pt] (2.0,0) -- (2.9,-0.4);
  % macula and fovea
  \fill[orange!60!brown!70] (0,0) circle (0.55);
  \fill[brown!70!black] (0,0) circle (0.14);
  \draw[black!60, dashed] (0,0) circle (0.95);
  % optic disc
  \fill[yellow!55!orange!60] (2.0,0.05) circle (0.42);
  \fill[yellow!25] (2.05,0.05) circle (0.18);
\end{scope}
\draw[black!70, thick] (0,0) circle (3.0);
% labels
\draw[-{Latex[length=1.6mm]}] (4.3,1.6) node[right, align=left] {Optic disc\\(vessels enter here)} -- (2.35,0.3);
\draw[-{Latex[length=1.6mm]}] (4.3,-0.4) node[right, align=left] {Superior and inferior\\vascular arcades} -- (1.2,-1.35);
\draw[-{Latex[length=1.6mm]}] (-4.1,1.3) node[left, align=right] {Macula\\(central vision)} -- (-0.75,0.55);
\draw[-{Latex[length=1.6mm]}] (-4.1,-0.8) node[left, align=right] {Fovea\\(sharpest vision)} -- (-0.12,-0.05);
\end{tikzpicture}}
\caption[Main anatomical landmarks of a colour fundus photograph]{Schematic of a colour
fundus photograph and its main landmarks. The optic disc is where the retinal vessels
enter the eye; the vascular arcades branch from it and surround the macula, whose centre,
the fovea, provides the sharpest vision. DR lesions near the macula threaten central vision
most directly.}
\label{fig:anatomy}
\end{figure}


\subsection{Global Burden}

Diabetes mellitus affects an estimated 537 million adults worldwide, a figure projected
to reach 783 million by 2045~\cite{idf2021atlas}. Its principal ocular complication,
diabetic retinopathy (DR), affects about 103 million people and is a leading cause of
preventable blindness in the working-age population~\cite{teo2021global,cheung2010dr}.
The disease is asymptomatic in its early, treatable stages, and timely laser or
anti-VEGF treatment greatly reduces the risk of severe visual loss~\cite{cheung2010dr},
provided that it is applied before tractional retinal detachment or extensive
neovascularisation.

\subsection{Pathophysiology and Stages of the Disease}
\label{sec:patho}

Chronic hyperglycaemia triggers polyol-pathway activation, protein kinase C
up-regulation, oxidative stress and the accumulation of advanced glycation end-products.
These converge on pericyte apoptosis, basement-membrane thickening and breakdown of the
blood--retinal barrier~\cite{cheung2010dr}. The healthy retinal vascular tree is
self-similar, with a fractal dimension close to $1.7$~\cite{masters2004fractal}, and its
complexity changes as the disease progresses. This change of vascular structure is the
biological motivation for the persistent-homology features studied in \cref{ch:toporet}.

Clinically, the cascade follows a predictable anatomical sequence rather than a uniform
deterioration. Capillary basement-membrane thickening and pericyte loss first produce
focal microaneurysms (Grade~1), localised weaknesses that later rupture or leak to form
haemorrhages and lipid-rich hard exudates (Grade~2). As capillary non-perfusion spreads,
the retina becomes ischaemic, producing venous calibre changes (beading) and
intraretinal microvascular abnormalities (Grade~3) -- compensatory shunt vessels that
bypass occluded capillary beds. Once ischaemia is extensive, VEGF up-regulation drives
neovascularisation at the optic disc or on the retinal surface (Grade~4), the
proliferative stage at which vitreous haemorrhage and tractional detachment become
acute risks.

This sequence is the clinical justification for two design choices of this thesis:
treating DR grading as an inherently \emph{ordinal} problem (\cref{sec:coral}) rather
than as a set of unrelated categories, and expecting the \emph{topological} complexity of
the vascular network to change monotonically with severity (\cref{sec:ph}).

\subsection{Clinical Signs of Diabetic Retinopathy}
\label{sec:signs}

The sequence described above produces a set of lesions that graders look for on
fundus photographs~\cite{wong2016dr,antonetti2012dr}:
\begin{description}[style=nextline,leftmargin=1.2em]
\item[Microaneurysms] small red dots, less than about $125~\mu$m, caused by out-pouching
of weakened capillary walls; they are the earliest visible sign of DR.
\item[Haemorrhages] red spots or blots caused by the rupture of capillaries; flame-shaped
haemorrhages lie in the nerve-fibre layer, dot-and-blot haemorrhages deeper in the retina.
\item[Hard exudates] bright yellow deposits of lipids and proteins that leak from damaged
vessels, often arranged in rings around a leaking area.
\item[Cotton-wool spots] fluffy white patches caused by small infarcts of the nerve-fibre
layer, a sign of local ischaemia.
\item[Venous beading and IRMA] irregular calibre of the veins and abnormal shunt vessels
inside the retina, both markers of severe ischaemia.
\item[Neovascularisation] new, fragile vessels growing on the optic disc or elsewhere on the
retina; they define proliferative DR and can bleed into the vitreous.
\end{description}
Diabetic macular oedema, the accumulation of fluid in the macula, can occur at any stage
of DR and is a major cause of visual loss. Its reliable diagnosis requires optical
coherence tomography (\cref{sec:imaging}); colour photographs show it only indirectly,
through exudates close to the fovea. This thesis focuses on the severity grade of DR
itself, as defined by the ICDR scale.

\subsection{Treatment and the Value of Early Detection}

Two treatments dominate the management of sight-threatening DR. Laser photocoagulation,
introduced in the 1970s and 1980s, destroys ischaemic peripheral retina to reduce the
drive for new vessels, and focal laser treats leaking areas near the macula.
Intravitreal injections of anti-VEGF agents, widely used since the 2010s, block the
growth factor responsible for neovascularisation and macular oedema. Vitrectomy is
reserved for advanced cases with vitreous haemorrhage or tractional detachment. All these
treatments are most effective when applied before severe visual loss; none of them can
restore tissue that is already destroyed~\cite{cheung2010dr,wong2016dr}. The clinical
value of screening therefore lies in detecting the disease \emph{before} symptoms appear
and in referring the right patients quickly, which is exactly what an automated grading
system must support.

\section{Retinal Imaging}
\label{sec:imaging}

\subsection{Colour Fundus Photography}

Colour fundus photography is the reference modality for DR screening. A fundus camera
illuminates the retina through the pupil and records a colour image of about 45 degrees
of the posterior pole, centred on the macula or on the optic disc. Images can be taken
with or without pharmacological dilation of the pupil (mydriatic or non-mydriatic
imaging); non-mydriatic cameras simplify screening but produce more ungradable images in
older patients or patients with cataract. Screening protocols use one or several fields
per eye, and the images are transmitted to readers or, increasingly, to automated
systems~\cite{abramoff2010retinal}.

\subsection{Other Imaging Modalities}

Optical coherence tomography (OCT)~\cite{huang1991oct} produces cross-sectional images
of the retinal layers with micrometre resolution and is the reference for detecting and
monitoring macular oedema. Fluorescein angiography, which requires an intravenous dye,
shows leakage and capillary non-perfusion but is invasive and reserved for specialist
care. Ultra-widefield imaging captures up to 200 degrees of the retina and reveals
peripheral lesions missed by standard fields. These modalities are complementary, but
colour fundus photography remains the only one that is cheap, fast and widespread enough
for population screening, and all the public grading datasets used in this thesis consist
of colour fundus photographs.

\subsection{Image Quality and Variability}

Fundus photographs vary widely across sites and devices: camera model, field of view,
resolution, illumination, compression and patient factors such as pupil size or media
opacities all change the appearance of the image. A non-negligible fraction of screening
images is ungradable. These variations are a major source of the performance drop observed
when a model trained in one population is applied to another, and they motivate two
design choices of this thesis: contrast normalisation with Contrast-Limited Adaptive
Histogram Equalisation (CLAHE) before any analysis, and evaluation on public datasets
acquired in three countries (United States, India and France), including external
datasets never seen during training.

\section{Severity Grading Scales}
\label{sec:grading}

The reference classification of DR severity was established by the Early Treatment
Diabetic Retinopathy Study (ETDRS), which defined a detailed scale of more than ten levels
from seven-field stereoscopic photographs~\cite{etdrs1991grading}. The ETDRS scale remains
the standard in clinical trials, but it is too detailed and too costly for routine
screening. The International Clinical Diabetic Retinopathy (ICDR) scale was proposed in
2003 to simplify it into five levels that can be graded from standard
photographs and that map directly onto clinical decisions~\cite{wilkinson2003icdr}.
Its grades are defined by the type of lesions present rather than by exact counts, with one
exception: the 4-2-1 rule, which separates severe from moderate NPDR by counting
haemorrhages, venous beading and IRMA per quadrant. The boundary between Grades~2 and~3,
decisive for referral, is therefore the only one defined by explicit counts.

\Cref{tab:icdr} summarises the ICDR scale used to grade every image in this thesis.
Grades~3 and~4 require urgent referral; missing them, that is, predicting No~DR or Mild
for such an image, is the error counted by the Severe-DR Non-Referral Rate (Sev-NR) for
Grade~3 and by the PDR Non-Referral Rate (PDR-NR) for Grade~4.

Grading the colour fundus photographs of \cref{sec:imaging} is a visual task:
the grader counts and locates microaneurysms, haemorrhages, exudates, venous beading and
new vessels, and maps them to a grade using the rules of \cref{tab:icdr}.

\begin{table}[H]
\centering
\caption[International Clinical Diabetic Retinopathy (ICDR) severity scale]{International Clinical Diabetic Retinopathy (ICDR) severity
scale~\cite{wilkinson2003icdr}. Grades 3--4 are safety-critical. The referral intervals
are indicative: they are not part of the ICDR definition and vary between screening
programmes.}
\label{tab:icdr}
\footnotesize
\begin{tabular}{@{}c l P{6.4cm} l@{}}
\toprule
Grade & Clinical name & Primary vascular signs & Referral urgency \\
\midrule
0 & No DR & None & Annual recall \\
1 & Mild NPDR & Microaneurysms only & 12 months \\
2 & Moderate NPDR & More than microaneurysms only (e.g.\ haemorrhages, exudates), below the 4-2-1 rule & 6 months \\
3 & Severe NPDR & 4-2-1 rule: $>20$ intraretinal haemorrhages in each of four quadrants, or venous beading in $\geq2$ quadrants, or IRMA in $\geq1$ quadrant; no PDR & 2--4 weeks \\
4 & Proliferative DR & Neovascularisation; vitreous or pre-retinal haemorrhage & Emergency \\
\bottomrule
\end{tabular}
\end{table}

\Cref{fig:fundus} shows one representative fundus photograph for each grade. The lesions of
\cref{tab:icdr} appear in increasing number and severity from left to right, but the
difference between two neighbouring grades is often a matter of a few lesions, which is the
source of the disagreement between graders discussed below.

\begin{figure}[H]
\centering
\setlength{\tabcolsep}{2pt}
\begin{tabular}{ccccc}
\includegraphics[width=0.188\textwidth]{fundus_grade0} &
\includegraphics[width=0.188\textwidth]{fundus_grade1} &
\includegraphics[width=0.188\textwidth]{fundus_grade2} &
\includegraphics[width=0.188\textwidth]{fundus_grade3} &
\includegraphics[width=0.188\textwidth]{fundus_grade4} \\[2pt]
\small\textbf{Grade 0} & \small\textbf{Grade 1} & \small\textbf{Grade 2} &
\small\textbf{Grade 3} & \small\textbf{Grade 4} \\
\footnotesize No DR & \footnotesize Mild NPDR & \footnotesize Moderate NPDR &
\footnotesize Severe NPDR & \footnotesize Proliferative DR
\end{tabular}
\caption[Representative fundus photographs for ICDR grades 0--4]{Representative fundus photographs for ICDR grades 0--4 (Kaggle EyePACS and
APTOS~2019 public datasets~\cite{kaggle2015eyepacs,aptos2019}). Lesion density and type
increase with grade: microaneurysms (Grade~1), haemorrhages and hard exudates (Grade~2),
venous beading and IRMA (Grade~3), neovascularisation (Grade~4).}
\label{fig:fundus}
\end{figure}

Grading is not an exact science. Even trained ophthalmologists disagree on a substantial
fraction of images, especially between adjacent grades, and the agreement depends on the
reference standard used: adjudication by a panel of retina specialists gives more
reliable labels than a single grader or a majority vote~\cite{krause2018grader}. Public
datasets are usually labelled by one or a few graders, so their labels contain some noise,
concentrated at the boundaries between grades. Any grading system, human or automated,
must therefore be judged against this intrinsic ambiguity. Many programmes grade each image
twice and send disagreements to a senior grader; an automated system is most useful when it
indicates which images deserve this second look. The calibrated probabilities and similar
reference cases of this thesis (\cref{ch:dots}) are first steps towards this role; an
automatic deferral of ambiguous images is left to future work. In the experiments of this thesis, the Mild--Moderate boundary concentrates the largest share of errors for every system (\cref{ch:toporet,ch:dots}), which is consistent with this intrinsic ambiguity.

\section{Screening Programmes and Patient Safety}
\label{sec:screening}

Population-scale screening is a public-health imperative, and automated
grading is widely seen as a way to make it more affordable at national
scale~\cite{ting2017jama,abramoff2018pivotal}. The safety of such a programme is measured by how many
sight-threatening cases it fails to refer. At a Severe-DR Non-Referral Rate of
$5.8\%$ -- the value of an EfficientNet-B3 baseline on the external Messidor-2 dataset in
this thesis -- a programme screening $500{,}000$ patients a year, of whom about $25{,}000$
have severe disease, would leave a projected $1{,}450$ severe cases unreferred every year (the assumptions of this
projection are stated in \cref{sec:impact}).

Organised screening has been shown to work at the scale of a country. In England, the
national programme started in 2003 invites every person with diabetes aged 12 and over for
annual photographic screening, with images graded in a multi-level quality-assured
process; since its introduction, diabetic retinopathy is no longer the leading cause of
certifiable blindness in the working-age population~\cite{scanlon2017english}. In the
United States, the EyePACS telemedicine network brought photographic screening to primary
care clinics by transmitting images to remote certified graders~\cite{cuadros2009eyepacs}
-- the same network that later provided the Kaggle dataset used in this thesis. Such
programmes need large numbers of trained graders, which is precisely the resource that
automated systems can save.

A screening programme is therefore judged on two different axes. Accuracy measures how
often the grade assigned to an image is correct, and is usually reported as the quadratic
weighted kappa (QWK) with the reference grader. Safety measures how often a patient who
needs urgent care is sent home, and is reported as the Severe-DR Non-Referral Rate
(Sev-NR). National programmes typically require a very low non-referral rate for
sight-threatening disease; this thesis adopts a target of $2\%$ throughout, used as a
predefined evaluation target rather than as a regulatory acceptance criterion. The two axes
do not move together: a system can gain several points of QWK while its Sev-NR stays above
the threshold, as the next section shows. The metrics are defined formally in
\cref{sec:metrics}.

\section{Public Datasets}
\label{sec:datasets}\label{app:datasets}

Research on automated DR grading relies on a small number of public datasets of colour
fundus photographs graded on the ICDR scale. Four of them are used in the DR experiments
of this thesis (\cref{tab:datasets}). Kaggle EyePACS~\cite{kaggle2015eyepacs} was released
for a public competition in 2015 and remains by far the largest. APTOS~2019~\cite{aptos2019}
was released for a second competition by an Indian tele-ophthalmology society.
Messidor-2~\cite{decenciere2014messidor} comes from French hospitals. KaggleDR-F1 is a
smaller Kaggle collection with five-grade labels, fused with APTOS~2019 in
\cref{ch:dots}. Other public datasets, such as IDRiD~\cite{porwal2018idrid,porwal2020idrid},
collected at a single Indian clinic with lesion segmentations, and DDR~\cite{li2019ddr},
from several Chinese hospitals, are also used in the literature; DRIVE~\cite{staal2004drive}
and STARE~\cite{hoover2000stare} provide manual vessel segmentations rather than grades and
have been central to the development of vessel segmentation methods.

\begin{table}[!htbp]
\centering
\caption[Grade distribution (\%) of the public datasets used in this thesis]{Grade distribution (\%) of the public datasets used in this thesis. Kaggle
EyePACS, APTOS~2019 and Messidor-2: counts distributed with the code of~\cite{belhadj2025plos}.}
\label{tab:datasets}
\small
\begin{tabular}{@{}lrccccc l@{}}
\toprule
Dataset & $N$ & G0 & G1 & G2 & G3 & G4 & Origin \\
\midrule
Kaggle EyePACS~\cite{kaggle2015eyepacs} & 35{,}126 & 73.5 & 7.0 & 15.1 & 2.5 & 2.0 & USA, multi-site \\
APTOS 2019~\cite{aptos2019} & 3{,}662 & 49.3 & 10.1 & 27.3 & 5.3 & 8.1 & India \\
Messidor-2~\cite{decenciere2014messidor} & 1{,}748 & 52.5 & 15.4 & 19.9 & 8.2 & 3.9 & France \\
KaggleDR-F1 & 2{,}197 & 49.3 & 10.1 & 27.3 & 5.3 & 8.1 & Kaggle collection \\
\bottomrule
\end{tabular}
\tabnote{$N$ counts retinal images, not independent patients: Kaggle EyePACS contains
both eyes of each patient. The full Kaggle EyePACS release has $88{,}702$ images; the
$35{,}126$ images with public labels are used in this thesis, as the TOPORET benchmark
(\enquote{Kaggle DR}) and as the external set of DGTS (\enquote{EyePACS}); $1{,}744$ of the $1{,}748$ Messidor-2 images are gradable.}
\end{table}

The datasets differ in ways relevant to the results of this thesis. Kaggle EyePACS, by
far the largest, aggregates images from many US clinical sites and camera models; it is
the most heterogeneous dataset and serves both as a TOPORET benchmark and as the external
validation set of DGTS. APTOS~2019 comes from an Indian screening programme. Messidor-2
(France, Topcon cameras) is a TOPORET benchmark and the external validation set of DOTS.
KaggleDR-F1 has exactly $60\%$ of the APTOS~2019 count in every grade, so its distribution
is identical to that of APTOS~2019; whether the two collections share images has not been
verified, a point discussed in \cref{ch:dots}. No dataset provides per-image demographic
annotations, so geographic diversity is only a proxy for population diversity. Grades~3
and~4 together are a minority in every dataset (from $4.5\%$ to $13.4\%$ combined), which
motivates class-balanced training: a system trained without correcting this imbalance would
under-predict exactly the two safety-critical grades on which Sev-NR and PDR-NR are
measured.

\section{Automated Grading: A Review of Methods}
\label{sec:decade}

Automated DR detection has been studied for more than two decades. Its history can be read
in four stages, which together explain the design of the systems developed in this
thesis. Recent reviews give a broader picture~\cite{tsiknakis2021review,alyoubi2020review,
ting2019ophthalmology}. The four stages are presented in chronological order, each with its main strength and the limitation that motivated the next one.

\subsection{Classical Image Analysis}

The first systems followed the reasoning of a human grader: detect the anatomical
landmarks, segment the vessels, find the lesions, and combine their counts into a
decision. Vessel segmentation received particular attention, with matched filters,
morphological operators and supervised pixel classifiers evaluated on the DRIVE and STARE
datasets~\cite{hoover2000stare,staal2004drive,fraz2012vessel}. Lesion detectors were
designed for bright lesions such as exudates~\cite{walter2002exudates} and for red lesions
such as microaneurysms and haemorrhages~\cite{niemeijer2005redlesions}. Combined into
complete systems, these methods reached screening-level sensitivity on large
populations~\cite{abramoff2008diabetes,abramoff2010retinal}. Their strength was
interpretability: every decision could be traced to detected lesions. Their weakness was
the accumulation of hand-tuned stages, each sensitive to image quality.

\subsection{Deep Convolutional Networks}

Deep learning replaced this pipeline by a single network trained end to end on graded
images~\cite{lecun2015deep,litjens2017survey}. In 2016, a convolutional network trained on
more than 100{,}000 graded images detected referable DR with a sensitivity and
specificity comparable to those of ophthalmologists~\cite{gulshan2016jama}, and similar
results were reported by other groups~\cite{abramoff2016improved,gargeya2017automated,
ting2017jama}. The Kaggle competition of 2015 popularised the quadratic weighted kappa as
the grading metric and led to many architectures based on standard ImageNet
backbones~\cite{he2016resnet,szegedy2016inception,tan2019efficientnet}. These networks
learn lesion detectors implicitly and are much more robust than hand-crafted pipelines,
but they are opaque, and their decisions are made one image at a time.

\subsection{Lesion-Aware and Attention Models}

A second generation of models tried to recover part of the interpretability of classical
systems. Lesion-mining approaches locate suspicious regions and use them to refine the
grade~\cite{quellec2017deep,wang2017zoomin}, while attention mechanisms weight the image
regions or the feature channels that matter most for each
grade~\cite{hu2018senet,he2021cabnet}. Systems trained on several tasks, such as lesion
segmentation and grading together, have extended detection to the whole disease spectrum,
including early stages~\cite{dai2021deepdr}. Visual explanations such as class activation
maps~\cite{zhou2016cam} or integrated gradients~\cite{sundararajan2017ig} have also been
shown to help human graders when displayed alongside the prediction~\cite{sayres2019integrated}.

\subsection{Transformers and Foundation Models}

Vision transformers~\cite{dosovitskiy2021vit} process an image as a sequence of patches
with self-attention, and capture long-range relations that convolutions see only
indirectly. Their main impact on DR grading has come through \emph{foundation models}:
networks pre-trained on very large unlabelled datasets and adapted to many tasks with
little labelled data~\cite{bommasani2021foundation}. RETFound~\cite{zhou2023retfound},
pre-trained by masked autoencoding on $1.6$ million retinal images, is the most prominent
example in ophthalmology and was reported to improve grading accuracy on several
benchmarks.

\subsection{Graph-Based and Population Approaches}

Graph neural networks, which learn on sets of related objects rather than isolated
samples, have been applied to many medical problems~\cite{ahmedt2021graphmed}. In
population graphs, each node is a patient and edges connect similar
patients~\cite{parisot2018disease}; the grade of a patient can then be informed by the
patients who resemble it. For DR, this line of work has remained limited to visual
similarity between images, and the global structure of the vessel network has not been
used to connect patients. \Cref{ch:background} introduces these methods in detail.

\subsection{Clinical Validation and Deployment}

Several systems have now been evaluated outside the laboratory. An autonomous system was
cleared by the US Food and Drug Administration in 2018 after a prospective trial in primary
care~\cite{abramoff2018pivotal}. Deep-learning systems were validated in national
programmes in India~\cite{gulshan2019india} and Thailand~\cite{ruamviboonsuk2019thailand},
and in a screening population in Zambia~\cite{bellemo2019zambia}. These studies confirm
that automated grading can match human graders on high-quality images, but they also
reveal practical obstacles: a study of a system deployed in Thai clinics found that
ungradable images, connectivity problems and workflow constraints reduced its usefulness
in practice~\cite{beede2020human}. More generally, the gap between retrospective accuracy
and clinical impact is now recognised as a central challenge of medical
AI~\cite{kelly2019challenges,liu2019comparison,rajpurkar2022ai}.

\subsection{Synthesis}

\Cref{tab:sota} summarises the systems most relevant to this thesis, from 2016 to 2026.
Accuracy has improved steadily, through better backbones, attention mechanisms and
retinal pre-training, but the three structural limitations defined in the next section
are not addressed together by any of the representative systems reviewed here.

\begin{table}[!htbp]
\centering
\caption[Comparison of DR grading systems, 2016--2026]{Comparison of DR grading systems,
2016--2026, with respect to the limitations L1--L3 of \cref{sec:limits}.}
\label{tab:sota}
\small
\setlength{\tabcolsep}{4pt}
\begin{tabular}{@{}l c P{5.2cm} c c c c@{}}
\toprule
System & Year & Key method & QWK & L1 & L2 & L3 \\
\midrule
Gulshan et al.~\cite{gulshan2016jama} & 2016 & Deep CNN (Inception-v3) & --$^{\dagger}$ & \xmark & \xmark & \xmark \\
CABNet~\cite{he2021cabnet} & 2021 & Category attention on a CNN backbone & --$^{\dagger}$ & \xmark & \xmark & \xmark \\
RETFound~\cite{zhou2023retfound} & 2023 & ViT-L/16, retinal pre-training & --$^{\dagger}$ & \xmark & \xmark & \xmark \\
SAMF-GB~\cite{belhadj2026samfgb}$^{\star}$ & 2026 & LightGBM $+$ structural features & 0.882 & \xmark & partial & \xmark \\
\midrule
TOPORET~\cite{belhadj2025plos}$^{\star}$ & 2026 & EfficientNet-B3 $+$ vessel TDA $+$ GraphSAGE & 0.88 & \cmark & \cmark & \xmark \\
DGTS~\cite{belhadj2026dgts}$^{\star}$ & 2026 & EfficientNet-B4 $+$ GraphSAGE $+$ PH regulariser & 0.941 & \cmark & partial & partial \\
DOTS~\cite{belhadj2026dots}$^{\star}$ & 2026 & EfficientNet-B3 $+$ FRE $+$ CORAL $+$ TTA & \textbf{0.953} & \xmark & \xmark & \cmark \\
\bottomrule
\end{tabular}
\tabnote{$^{\dagger}$Not reported in a form comparable to five-grade QWK (for example,
Gulshan et al.\ report AUC $=0.991$ for referable DR). $^{\star}$This thesis; QWK on the
dataset and protocol of the corresponding chapter (TOPORET: Kaggle DR, transductive five-fold cross-validation; DGTS and DOTS:
five-fold cross-validation on the fused APTOS~2019 $+$ KaggleDR-F1 corpus), not directly
comparable across rows. ``Partial'': DGTS uses feature-space rather than vascular topology,
and reports but does not meet the safety criterion on external data.}
\end{table}

Two patterns emerge from \cref{tab:sota}.

\paragraph{Pattern 1: accuracy gains do not resolve structural limitations.}
Better backbones and retinal pre-training improve accuracy, yet L1, L2 and L3 remain
unaddressed in each of the reviewed external systems. This suggests that the limitations
relate to architectural choices rather than to hyperparameters or dataset size, although
indirect mechanisms (for example global attention in transformers) may capture part of
the missing information.

\paragraph{Pattern 2: safety is rarely measured.}
None of the reviewed external systems reports the rate at which severe or proliferative
cases would be left unreferred under a five-grade protocol, and few are validated on a
dataset acquired in another country. A high QWK does not guarantee a low non-referral
rate: in \cref{ch:dots}, DGTS generalises well in QWK but its Sev-NR on external data
remains slightly above $2\%$.

\section{Challenges and Limits of Current Systems}
\label{sec:challenges}

Despite a decade of progress, automated DR grading still faces open challenges. Some are
practical: models trained on one population lose accuracy on images from other cameras or
other countries; public datasets are small for the severe grades; and a model that always
answers, even when the image is ambiguous, is hard to accept in a clinical workflow. This
thesis focuses on three limitations that are not practical but \emph{architectural}: they
recur across the representative systems reviewed above and are unlikely to be removed by
more data or larger networks alone. They are defined below as properties of an
architecture, so that each system of \cref{tab:sota} can be checked against them.

\subsection{Three Structural Limitations}
\label{sec:limits}

The following limitations are defined as binary architectural properties. They do not
depend on the volume of training data, on network depth or on computational budget.

\paragraph{Definition 1.1 (L1 -- No Relational Context).}
A system satisfies L1 if and only if $P(Y_i \mid I_i, \mathcal{C}) = P(Y_i \mid I_i)$,
where $\mathcal{C}$ denotes the patient cohort. All the per-image systems of
\cref{tab:sota} satisfy this condition.

\paragraph{Definition 1.2 (L2 -- No Vascular Topology).}
A system satisfies L2 if and only if its feature map $\varphi : I \to \mathbb{R}^d$ uses
only operations with bounded receptive fields (convolutions, local attention). Such a
map does not explicitly compute global topological invariants of the vascular network.

\paragraph{Definition 1.3 (L3 -- No Safety-Oriented Grading).}
A system satisfies L3 if and only if it is trained with an objective that ignores the
order of the grades and is evaluated without measuring the non-referral of severe and
proliferative cases against a safety threshold. Such a system cannot show that its errors
fall on the safe side of the referral decision.

\subsection{Why the Limitations Are Studied Together}
\label{sec:coupled}

A natural question is whether each limitation could be resolved by adding one component
to an existing system. The experiments of this thesis suggest that the answers interact.

Resolving \textbf{L1 alone} (adding a population graph built from visual similarity only,
as in DGTS) improves QWK by $0.018$ (\cref{tab:dgts-abl}), but leaves the topology of the
vessels unused; a dual metric that includes a topological distance requires the
topological pipeline of L2 to be already in place.

Resolving \textbf{L2 alone} (adding topological features without the graph) yields $+0.02$
QWK (\cref{tab:ablation3}) -- useful, but partial. The largest gain ($+0.04$ QWK) is
obtained only when topology is used both as a patient feature and inside the similarity of
the graph of L1, and the same complementarity is observed with DGTS under an inductive protocol
(\cref{ch:dots}).

Resolving \textbf{L3} requires more than an ordinal loss: on Messidor-2, CORAL alone leaves
Sev-NR at $3.7\%$, and only the complete DOTS pipeline reaches $1.4\%$ (\cref{tab:h3b}).
Conversely, DGTS improves accuracy and generalisation through L1 and L2 but does not meet
the safety threshold on external data.

\looseness=-1
This interaction is the motivating rationale for studying the three limitations in one
research programme (\cref{ch:toporet,ch:dots}). The thesis addresses them in complementary
systems; combining all three in one model is left to future work.

\subsection{Positioning of This Thesis}
\label{sec:position}

\Cref{tab:sota} makes the gap concrete: over a decade of progress in accuracy, none of the
representative external systems resolves the three structural limitations, because
per-image classifiers neither relate patients, nor compute global topological invariants,
nor target the referral decision (\cref{sec:limits}). This thesis therefore investigates
architectural additions rather than further scaling; \cref{tab:limits} summarises the three
limitations, their clinical consequence and the component that addresses each of them.

\begin{table}[!htbp]
\centering
\caption[The three structural limitations and how this thesis addresses them]{The three structural limitations and how this thesis addresses them.}
\label{tab:limits}
\footnotesize
\begin{tabular}{@{}P{3.0cm} P{4.6cm} P{5.6cm} c@{}}
\toprule
Limitation & Clinical consequence & Component of this thesis & Ch. \\
\midrule
L1 -- no relational context & Similar patients never inform each other's grade & Population graphs (TOPORET; inductive in DGTS) & 3--5 \\
L2 -- no vascular topology & Loss of vessel loops and rarefaction is not measured & Persistent homology of the vessel skeleton (TOPORET); topological regulariser (DGTS) & 3--5 \\
L3 -- no safety-oriented grading & Severe cases may be graded as No DR or Mild, unmeasured & CORAL ordinal head and safety evaluation with external validation (DOTS) & 5 \\
\bottomrule
\end{tabular}
\end{table}

\section*{Conclusion}
\phantomsection\addcontentsline{toc}{section}{Conclusion}

\looseness=-1
Despite a decade of gains in accuracy, the representative systems reviewed in this chapter
still grade each image in isolation (L1), ignore the topology of the vascular tree (L2) and
are rarely trained or evaluated with the referral of severe cases in mind (L3). Addressing
them requires learning on graphs, topological data analysis, ordinal prediction and
calibrated safety evaluation, which the next chapter introduces with the evaluation metrics.
